\section{Problem Formulation and Scope}
\label{sec:problem}

\subsection{Notation and Representation of Legal Rules}

Let \(\mathcal{R}=\{r_1,r_2,\ldots,r_n\}\) denote the set of legal rules
normalized from the corpus. Each rule \(r_i\) is represented by the tuple:

\begin{equation}
r_i =
\left(
    c_i,\,
    e_i,\,
    s_i,\,
    a_i
\right),
\label{eq:rule}
\end{equation}

where:

\begin{itemize}
    \item \(c_i\) is the applicability condition of the rule;
    \item \(e_i\) is the legal effect that arises when the condition is
    satisfied;
    \item \(s_i\) is the relevant legal subject;
    \item \(a_i\) is the statutory basis of the rule.
\end{itemize}

Condition \(c_i\) and effect \(e_i\) are mapped to the corresponding legal
events \(v_i^{c}\) and \(v_i^{e}\). Within the scope of the model, rule
\(r_i\) defines a directed relation:

\begin{equation}
v_i^{c}
\xrightarrow{r_i}
v_i^{e}.
\end{equation}

This relation is referred to as a legal condition--effect relation. The term
\emph{causal} is used in this paper in this operationalized sense: the
condition event activates a rule mechanism, and the rule mechanism gives
rise to the effect event. The relation is not interpreted as a causal
relation discovered or estimated from observational data.

The corpus version used in the experiments currently maps each rule
primarily to one positive condition and one positive effect. More complex
structures, such as multiple simultaneous conditions, exceptions, negation,
priority orderings among rules, or normative conflicts, are not yet fully
represented. This limitation is discussed in detail in
Section~\ref{sec:limitations}.

\subsection{Legal Causal Graph}

From the rule set \(\mathcal{R}\), the system constructs a heterogeneous
graph:

\begin{equation}
\mathcal{G}
=
\left(
    \mathcal{V},
    \mathcal{E}
\right),
\end{equation}

with the node set:

\begin{equation}
\mathcal{V}
=
\mathcal{V}_{\mathrm{event}}
\cup
\mathcal{V}_{\mathrm{rule}}
\cup
\mathcal{V}_{\mathrm{article}},
\end{equation}

where:

\begin{itemize}
    \item \(\mathcal{V}_{\mathrm{event}}\) contains legal events;
    \item \(\mathcal{V}_{\mathrm{rule}}\) contains normalized rules;
    \item \(\mathcal{V}_{\mathrm{article}}\) contains statutory provisions
    or legal authorities.
\end{itemize}

The edge set \(\mathcal{E}\) represents relations among events, rules, and
statutory provisions. In particular, \texttt{CAUSES} edges connect condition
events to effect events. A pair of events may be supported by multiple
different rules; the graph therefore permits multiple edges or multiple rule
mechanisms between the same pair of nodes.

Storing events, rules, and statutory provisions together enables the system
to perform three related but non-identical tasks:

\begin{enumerate}
    \item retrieving events relevant to the query;
    \item recovering causal paths among events;
    \item tracing each path hop back to its supporting rule and statutory
    provision.
\end{enumerate}

\subsection{Causal Paths and Multi-Hop Reasoning}

A causal path of length \(h\) is defined as an ordered chain:

\begin{equation}
p =
\left(
    v_0,\,
    r_1,\,
    v_1,\,
    r_2,\,
    \ldots,\,
    r_h,\,
    v_h
\right),
\label{eq:path}
\end{equation}

such that for every \(j \in \{1,\ldots,h\}\), rule \(r_j\) supports the
relation:

\[
v_{j-1}
\xrightarrow{r_j}
v_j.
\]

Here:

\begin{itemize}
    \item \(v_0\) is the seed event or initial cause;
    \item \(v_h\) is the outcome event;
    \item \(v_1,\ldots,v_{h-1}\) are mediators;
    \item \(r_1,\ldots,r_h\) are the rules supporting the individual hops.
\end{itemize}

In the main experimental configuration, the system retrieves paths of at
most two hops, that is, \(h \leq 2\). A two-hop path has the form:

\[
v_0
\xrightarrow{r_1}
v_1
\xrightarrow{r_2}
v_2,
\]

where \(v_1\) is both the effect of rule \(r_1\) and the condition of rule
\(r_2\).

The two-hop limit was selected because the current benchmark is constructed
primarily from two-hop chains present in the graph. This is not a
theoretical limitation of the path representation; however, increasing the
number of hops substantially expands the candidate space and the potential
for retrieval errors to accumulate.

\subsection{Distinguishing Direct and Indirect Relations}

For two events \(A\) and \(C\), a direct relation exists when the graph
contains at least one rule mechanism:

\[
A \xrightarrow{r} C.
\]

Conversely, if the graph contains only a path:

\[
A \xrightarrow{r_1} B \xrightarrow{r_2} C,
\]

then the system may conclude that \(A\) is related to \(C\) through mediator
\(B\), but it must not conclude that a direct edge \(A \rightarrow C\)
exists.

This distinction is particularly important for questions such as:

\begin{quote}
``Does \(A\) directly lead to \(C\)?''
\end{quote}

If the graph supports only a two-hop path through \(B\), the direct claim
must be rejected even if \(C\) can still be inferred indirectly from \(A\).

\subsection{Queries and Claim Types}

Let \(q\) denote a natural-language query. The system analyzes \(q\) to
identify the type of relation requested by the user. The main claim types
are:

\begin{description}
    \item[\texttt{STANDARD\_CAUSAL}] The query requests the identification or
    explanation of a standard condition--effect relation.

    \item[\texttt{DIRECT\_CAUSAL\_CLAIM}] The query asserts or asks about the
    existence of a direct edge between two events.

    \item[\texttt{REMOVE\_MEDIATOR\_OUTCOME\_REMAINS}] The query claims that
    the outcome still occurs when the mediator is removed.

    \item[\texttt{REMOVE\_MEDIATOR\_OUTCOME\_DISAPPEARS}] The query claims
    that the outcome disappears when the mediator is removed.

    \item[\texttt{MEDIATOR\_NECESSARY\_CLAIM}] The query asks whether the
    mediator is necessary for the outcome.

    \item[\texttt{COUNTERFACTUAL\_UNSPECIFIED}] The query contains a
    counterfactual signal but does not specify whether the outcome is
    expected to remain or disappear.
\end{description}

Claim identification relies only on the query content and the retrieved
event structure. The system does not use \texttt{question\_type}, the
reference decision label, the gold path, or the gold answer from the
benchmark.

\subsection{Factual World and Counterfactual World}

Each event \(v \in \mathcal{V}_{\mathrm{event}}\) is associated with a state
variable:

\begin{equation}
X_v
\in
\left\{
    \mathrm{TRUE},
    \mathrm{FALSE},
    \mathrm{UNKNOWN}
\right\}.
\end{equation}

The three states are interpreted as follows:

\begin{itemize}
    \item \(\mathrm{TRUE}\): the event is inferred to occur in the world
    under consideration;
    \item \(\mathrm{FALSE}\): the event is inferred not to occur;
    \item \(\mathrm{UNKNOWN}\): the evidence is insufficient to determine
    the event state, or the reasoning mechanisms conflict.
\end{itemize}

The factual world \(W^{F}\) is determined from the seed event and the rule
mechanisms activated in the graph. For a causal path:

\[
A \rightarrow M \rightarrow Y,
\]

the system first checks the factual states:

\[
X_A^{F},\quad X_M^{F},\quad X_Y^{F}.
\]

To evaluate the role of mediator \(M\), LegalSCM performs the hard
intervention:

\begin{equation}
do
\left(
    X_M = \mathrm{FALSE}
\right).
\label{eq:do}
\end{equation}

Under the intervention semantics of SCMs
\cite{pearl2009causality}, the operation in Equation~\ref{eq:do} does not
merely change the value of \(M\). It also replaces the ordinary structural
mechanism that generates \(M\), equivalently severing the mediator's
incoming mechanisms in the intervened world.

The corresponding counterfactual world is denoted by:

\[
W^{CF}_{M}
=
W^{F}
\mid
do
\left(
    X_M = \mathrm{FALSE}
\right).
\]

After the intervention, LegalSCM re-infers the outcome state:

\[
X_Y^{CF}
=
X_Y
\left(
    W^{CF}_{M}
\right).
\]

Comparing \(X_Y^{F}\) with \(X_Y^{CF}\) distinguishes the following cases:

\begin{itemize}
    \item if
    \(X_Y^{F}=\mathrm{TRUE}\) and
    \(X_Y^{CF}=\mathrm{FALSE}\),
    the mediator is considered necessary within the model under
    consideration;

    \item if
    \(X_Y^{F}=\mathrm{TRUE}\) and
    \(X_Y^{CF}=\mathrm{TRUE}\),
    the outcome can still be sustained after the intervention;

    \item if either state is
    \(\mathrm{UNKNOWN}\),
    the intervention result is insufficient for a definitive conclusion.
\end{itemize}

The notion of ``necessity'' here is meaningful only relative to the rule
mechanisms and factual context supplied to LegalSCM. It does not establish
legal necessity in every real-world situation.

\subsection{Path Ablation}

As a baseline, the study uses a path-ablation procedure. For the same path:

\[
A \rightarrow M \rightarrow Y,
\]

mediator \(M\) is removed from the graph, after which the system searches for
alternative paths from \(A\) to \(Y\). If no suitable path remains, the
mediator is considered necessary according to the reachability criterion. If
an alternative path reaches the score threshold, the outcome is considered
still reachable.

Path ablation and hard intervention address two different questions:

\begin{itemize}
    \item path ablation asks whether the graph contains a path from \(A\) to
    \(Y\) after \(M\) is deleted;

    \item LegalSCM asks how the state of \(Y\) changes when \(M\) is forcibly
    assigned \(\mathrm{FALSE}\) and the relevant mechanisms are re-inferred.
\end{itemize}

Accordingly, the path-ablation result is not interpreted as the result of a
\(do(\cdot)\) operation. In this study, path ablation is used only as a
baseline for evaluating differences in quality, cost, and trace detail.

\subsection{Verification Decision}

Given a query \(q\), a set of candidate paths \(\mathcal{P}_q\), and the
claim-analysis result, the system produces a decision:

\begin{equation}
D_q
\in
\left\{
    \texttt{SUPPORTED},
    \texttt{REJECT\_DIRECT\_CLAIM},
    \texttt{UNCERTAIN}
\right\}.
\end{equation}

The labels are interpreted as follows:

\begin{description}
    \item[\texttt{SUPPORTED}] The graph structure and verification result
    support the claim detected in the query.

    \item[\texttt{REJECT\_DIRECT\_CLAIM}] A direct claim, or another claim
    requiring rejection, is inconsistent with the path structure and
    verification result. For example, the query asserts that \(A\) directly
    leads to \(C\), but the graph supports only
    \(A \rightarrow B \rightarrow C\).

    \item[\texttt{UNCERTAIN}] The system cannot identify a sufficiently
    reliable path, the factual path cannot be confirmed, or the intervention
    result contains an undetermined state.
\end{description}

In addition to label \(D_q\), the system records a decision score, an
explanation, the primary path, the claim type, and the relevant factual and
counterfactual fields. These fields are passed unchanged to the answer
generation stage to reduce the risk that the answer generator produces a
conclusion that contradicts the verification stage.

\subsection{Answer Formats}

The claim type determines the semantics of the decision, whereas the answer
intent determines the form of information to be presented. The system
supports the following answer intents:

\begin{description}
    \item[\texttt{FINAL\_OUTCOME}] Return the final effect of the primary
    path.

    \item[\texttt{INITIAL\_CAUSE}] Return the initial cause or event.

    \item[\texttt{BRIDGE\_EVENT}] Return the mediator between the cause and
    the outcome.

    \item[\texttt{CHAIN\_EXPLANATION}] Present each hop in sequence together
    with its corresponding statutory basis.

    \item[\texttt{YES\_NO}] Return yes, no, or insufficient evidence, as
    appropriate for the verification decision.

    \item[\texttt{BRANCH}] Present multiple outcomes that share a common
    causal prefix.

    \item[\texttt{CONVERGENCE}] Present multiple initial conditions that
    converge on the same suffix or outcome.

    \item[\texttt{CLAIM\_VERIFICATION}] Answer directly whether a direct or
    counterfactual claim is true or false.
\end{description}

The answer intent is inferred from the query and the query-analysis result
from the verification stage. As with claim detection, this process does not
read a question-type label or gold answer.

\subsection{Problem Output}

The complete pipeline can be described as the mapping:

\begin{equation}
F:
\left(
    q,\,
    \mathcal{R},\,
    \mathcal{G}
\right)
\longrightarrow
\left(
    \widehat{\mathcal{E}}_q,\,
    \widehat{\mathcal{P}}_q,\,
    D_q,\,
    A_q,\,
    C_q,\,
    M_q
\right),
\end{equation}

where:

\begin{itemize}
    \item \(\widehat{\mathcal{E}}_q\) is the set of retrieved evidence;
    \item \(\widehat{\mathcal{P}}_q\) comprises the candidate causal paths
    and the selected primary path;
    \item \(D_q\) is the verification decision;
    \item \(A_q\) is the final answer;
    \item \(C_q\) is the set of citations used;
    \item \(M_q\) is metadata describing retrieval, the factual world, the
    intervention, the counterfactual world, and answer generation.
\end{itemize}

An output is considered successful when it simultaneously satisfies several
criteria:

\begin{enumerate}
    \item the relevant rules and events are retrieved;
    \item the primary path reflects the correct reasoning order;
    \item the verification decision is consistent with the claim type;
    \item the answer matches the answer intent;
    \item the citations genuinely support the hops used;
    \item the metadata makes the processing procedure traceable and
    reproducible.
\end{enumerate}

\subsection{Scope and Assumptions}

The study is conducted under the following assumptions and limitations:

\begin{enumerate}
    \item \textbf{Normative relations rather than causal discovery.}
    Edge directions are obtained from the condition--effect structure of
    legal rules and are not learned from observational data.

    \item \textbf{Deterministic three-valued reasoning.}
    LegalSCM does not represent probabilities. Events are assigned
    \(\mathrm{TRUE}\), \(\mathrm{FALSE}\), or \(\mathrm{UNKNOWN}\).

    \item \textbf{Limited interventions.}
    The current experiments focus only on setting a mediator to
    \(\mathrm{FALSE}\). Multivariable interventions and interventions on
    continuous values are outside the scope of this study.

    \item \textbf{Short paths.}
    The main evaluation configuration limits causal paths to at most two
    hops.

    \item \textbf{Simplified rule representation.}
    The current corpus does not fully cover exceptions, negation, multiple
    simultaneous conditions, or priority among rules.

    \item \textbf{Extractive answer generation.}
    The main experiments use a template-based answer generator and do not
    evaluate the expressive capability of a free-form generative language
    model.

    \item \textbf{Benchmark derived from the same graph.}
    The questions are generated from the graph used in the pipeline.
    Consequently, the evaluation results reflect the ability to recover and
    process encoded structures rather than fully demonstrating
    generalization to an independent legal corpus.

    \item \textbf{Not a substitute for legal advice.}
    The system is developed for research purposes and must not be used as a
    source of official legal advice or legal conclusions.
\end{enumerate}
